% !PRINTED none
\weberpage{none}
% [source note: this Abschnitt opener has no printed folio; sequence position corresponds to inferred page 670.]

\begin{center}
\gesperrt{Dreiundzwanzigster Abschnitt.}

\gesperrt{Classenzahl.}

\rule{4em}{0.4pt}

\S.\,185.

\gesperrt{Der Dirichlet'sche Satz über die Einheiten.}
\end{center}

Die Untersuchungen über die Anzahl der Idealclassen in einem algebraischen Körper, die uns nun die nothwendigen Ergänzungen der Beweise des vorigen Abschnittes geben sollen, bedienen sich der Methoden, die \gesperrt{Dirichlet} in die Zahlentheorie eingeführt hat\begingroup\renewcommand{\thefootnote}{\arabic{footnote})}\footnote{\gesperrt{Dirichlet's}, Recherches sur diverses applications de l'analyse infinitésimale à la théorie des nombres. \gesperrt{Crelle's} Journal, Bd. 19, 21. \gesperrt{Dirichlet's} Werke, Bd. I (1839, 1840). (Hierher gehören übrigens auch die nachgelassenen Abhandlungen von \gesperrt{Gauss}, \glqq De nexu inter multitudinem classium etc.\grqq{} \gesperrt{Gauss'} Werke, Bd. II.) Die Anwendung dieser Principien auf die Kreistheilungszahlen hat zuerst \gesperrt{Kummer} gemacht [\gesperrt{Crelle's} Journal, Bd. 35, 40 (1847, 1850); \gesperrt{Liouville's} Journal, Bd. 16 (1851)]. Die allgemeine Theorie der Einheiten ist gleichfalls von \gesperrt{Dirichlet} begründet [\gesperrt{Dirichlet's} Werke, Bd. I, S. 622, 633, 639 (1840, 1842, 1846)]. Die allgemeine Theorie der Einheiten und die Bestimmung der Classenzahlen hat \gesperrt{Dedekind} gegeben: \gesperrt{Dirichlet-Dedekind}, Vorlesungen über Zahlentheorie, 4. Auflage, \S.\,183 f. \gesperrt{Minkowski}, Geometrie der Zahlen.}\endgroup.

Sie sind nicht mehr rein algebraisch, insofern dabei auch transcendente Functionen und Zahlen, wie der natürliche Logarithmus und die Zahl $\pi$ vorkommen. Auch wird von Grenzübergängen, wie in der Integralrechnung, Gebrauch gemacht, aus denen wir schon in \S.\,156 ein schönes Resultat gezogen haben.

Die Untersuchungen, von denen wir zunächst zu handeln haben, sind ganz allgemein, d.\,h. sie gelten für jeden algebraischen Zahlkörper, und es bietet keinen Vortheil der Einfachheit, sie auf specielle Körper, wie etwa die Kreistheilungskörper, zu beschränken.
